\documentclass[11pt,reqno]{amsart}

\usepackage{amsmath,amssymb,amsthm}
\usepackage[margin=1.15in]{geometry}
\usepackage[colorlinks=true,linkcolor=blue,citecolor=blue,urlcolor=blue]{hyperref}

\theoremstyle{plain}
\newtheorem{theorem}{Theorem}[section]
\newtheorem{proposition}[theorem]{Proposition}
\newtheorem{lemma}[theorem]{Lemma}
\newtheorem{corollary}[theorem]{Corollary}

\newcommand{\opt}{\beta}
\newcommand{\Nb}{N}
\newcommand{\optf}{\beta^{*}}
\newcommand{\dens}{\rho}

\begin{document}

\title[Sparse halves in triangle-free graphs]
{Sparse halves in triangle-free graphs: a bound of $131/5000$}

\author{Amir Sarid}
% amsart prints \date at the foot of the first page.  Show it under the author
% instead.  \dedicatory is the one title-block hook between author and abstract.
\dedicatory{\normalfont August 4, 2026}
\date{}

\begin{abstract}
Erd\H{o}s conjectured that every triangle-free graph on $n$ vertices contains
$\lfloor n/2\rfloor$ vertices spanning at most $n^{2}/50$ edges, and offered a
prize for a proof.  The best bound previously known is $\tfrac{27}{1024}n^{2}$, proved
by Razborov using the flag-algebra method.  We prove that every triangle-free
graph on $n$ vertices contains $\lfloor n/2\rfloor$ vertices spanning at most
\[
 \frac{131}{5000}n^{2}=0.0262\,n^{2}
\]
edges.  The argument uses the max-cut theorem of Balogh, Clemen and
Lidick\'y, an exact rational flag-algebra bound on the four-cycle density, and
a perturbation supported on the neighbourhood of a single vertex.

The proof has been formalized in Lean, with the max-cut theorem taken as a
hypothesis in the form its authors state it, Theorem~2(a) of \cite{bcl}.
Everything else is proved inside the formal development.  OpenAI's
GPT-5.6~Sol and Anthropic's Claude Fable~5 assisted in the research,
writing and formalization of this work.
\end{abstract}

\maketitle

\section{Introduction}

For a graph $G$ on $n\ge1$ vertices let
\[
 \opt(G)=\frac1{n^2}\min\bigl\{e(S):S\subseteq V(G),\
 |S|=\lfloor n/2\rfloor\bigr\},
\]
where $e(S)$ is the number of edges spanned by $S$.  The normalization by
$n^{2}$ is the one used in the literature.  Erd\H{o}s \cite{erdos75} asked whether every triangle-free graph satisfies
$\opt(G)\le1/50$.  He returned to the question repeatedly, offering a prize of
\$250 for it in \cite{erdos97}, and Keevash and Sudakov describe it as one of
his favourite conjectures \cite{keevash-sudakov}.  See also
\cite{erdosproblems}.  Balanced blow-ups of $C_{5}$ attain $\opt=1/50$ whenever the blow-up factor is
even, and balanced blow-ups of the Petersen graph attain it for every factor,
so $1/50$ is the smallest constant for which such a bound can hold.  These are
the only extremal examples known \cite{razborov-halves}.

For a general triangle-free graph, the bound $\opt(G)\le1/16$ is immediate
from a random half.  Erd\H{o}s, Faudree, Rousseau and Schelp \cite{efrs}
proved $\opt(G)\le1/30$, and Krivelevich \cite{krivelevich} improved this to
$\opt(G)\le1/36$.  The conjecture itself is known in several ranges of the
edge density $\dens=\dens(G)=2e(G)/n^{2}$.  It is immediate from a random half
when $\dens\le4/25$.  Keevash and Sudakov \cite{keevash-sudakov} proved it for
$\dens\le1/6$ and for $\dens\ge2/5$, showing moreover that in the latter range
the balanced blow-up of $C_{5}$ is the only extremal example.  For regular
graphs the case $\dens\ge2/5$ was already known
\cite[Theorem~3]{krivelevich}.  Norin and Yepremyan
\cite{norin-yepremyan} relaxed $\dens\ge2/5$ to
$\dens\ge\tfrac25-\varepsilon_{0}$ for an absolute $\varepsilon_{0}>0$, and
proved the conjecture for graphs of normalized minimum degree at least
$5/14$.  The best
general bound previously known is due to Razborov \cite{razborov-halves}, who
proved
\[
 \opt(G)\le\frac{27}{1024}
\]
by the flag-algebra method \cite{razborov-flag}, and also settled the
conjecture for graphs of girth at least five, for graphs of independence
number at least $2n/5$, and for strongly regular graphs.

A companion question of Erd\H{o}s asks how many edges must be deleted from a
triangle-free graph to make it bipartite.  It was studied by Erd\H{o}s,
Faudree, Pach and Spencer \cite{efps} and by Erd\H{o}s, Gy\H{o}ri and
Simonovits \cite{egs}, and most recently by Balogh, Clemen and Lidick\'y
\cite{bcl}, whose theorem is the one result the present proof quotes.
There, one asks for a bipartition whose two sides span few edges, whereas here
a single set, of size exactly $\lfloor n/2\rfloor$, has to be sparse.

Our result is the following:

\begin{theorem}\label{thm:main}
Every finite triangle-free graph $G$ on $n\ge1$ vertices satisfies
\[
 \opt(G)\le\frac{131}{5000}=0.0262.
\]
\end{theorem}

We treat separately the graphs with $\dens\ge7/22$ and those with
$\dens<7/22$.

For $\dens\ge7/22$ a double counting argument over the edges suffices.  If
$\{u,v\}$ is an edge of a triangle-free graph, then $\Nb(u)$ and $\Nb(v)$ are
disjoint independent sets, so each of them is the sparse part of a candidate of
the required total size, and the cost of the better of the two candidates is
governed by the number of edges between $\Nb(u)$ and $\Nb(v)$.  Averaging over
all edges of $G$ turns that quantity into the number of homomorphisms of the
four-cycle into $G$, which we bound from below by an exact rational
flag-algebra certificate.  This construction goes back to Krivelevich
\cite{krivelevich}, and the averaged form we use is due to Razborov
\cite{razborov-halves}.

For $\dens<7/22$ we use the max-cut theorem of Balogh, Clemen and Lidick\'y
\cite{bcl}, which supplies a bipartition of $V(G)$ with at most $2n^{2}/47$
edges inside its two sides.  Two candidates are built from such a bipartition,
and each is improved by a perturbation supported on the neighbourhood of a
single vertex.  A harmonic-mean averaging argument produces
a vertex at which that improvement is cubic in the edge density of the larger
side.  What is left is to maximize an explicit function of two parameters.  Its
maximum occurs when the larger side of the bipartition is $417/241$ times the
smaller, and the value there is below $131/5000$.

Section~\ref{sec:background} fixes notation, relates $\opt$ to a fractional
relaxation, and records the blow-up invariance that removes the largeness
hypothesis from the max-cut theorem.  Section~\ref{sec:high} treats
$\dens\ge7/22$ and Section~\ref{sec:low} treats $\dens<7/22$.
Theorem~\ref{thm:main} is proved at the end of Section~\ref{sec:low}, and
Section~\ref{sec:verification} describes the verification.

The constant $131/5000$ is close to, but not exactly, the best this argument
gives.  The bound of Section~\ref{sec:high} decreases as the split density
grows while that of Section~\ref{sec:low} increases, so the two are in balance
at one value of $\dens$.  Numerically that value is near $0.3184$, where both
are about $0.026176$.  We split at the nearby short rational $\dens=7/22$ and
state a short decimal, which together cost about $2\cdot10^{-5}$.

\section{Fractional halves and graphons}\label{sec:background}

Define
\begin{equation}\label{eq:b-def}
 \optf(G)=\frac1{n^2}\min\left\{
 \sum_{\{x,y\}\in E(G)}f_xf_y:
 f:V(G)\to[0,1],\ \sum_x f_x=\frac n2
 \right\}.
\end{equation}

A function $f:V(G)\to[0,1]$ with $\sum_xf_x=n/2$ is a \emph{fractional half},
and $n^{-2}\sum_{\{x,y\}\in E(G)}f_xf_y$ is its \emph{cost}, so $\optf(G)$ is the
least cost of a fractional half.

\begin{lemma}\label{lem:rounding}
$\opt(G)\le\optf(G)$ for every graph $G$.
\end{lemma}

\begin{proof}
Choose a minimizer of \eqref{eq:b-def} with as few nonintegral coordinates as
possible.  If $f_x,f_y\notin\{0,1\}$, replace $(f_x,f_y)$ by
$(f_x+t,f_y-t)$.  On the feasible interval for $t$, the objective is affine
when $\{x,y\}\notin E(G)$ and is a concave quadratic when $\{x,y\}\in E(G)$.  One
endpoint is therefore no worse, contradicting the choice of $f$.  Thus at
most one coordinate is nonintegral.  For even $n$ there is none.  For odd
$n$ the remaining coordinate is $1/2$, and deleting its vertex gives
$\lfloor n/2\rfloor$ vertices without adding an edge.
\end{proof}

For a graph $H$ write
\[
 t(H,G)=\frac{\hom(H,G)}{n^{|V(H)|}}
\]
for the homomorphism density of $H$ in $G$.  Here $\hom(H,G)$ is the number of
maps $V(H)\to V(G)$ carrying every edge to an edge, not necessarily injective.
Thus $\dens=t(K_2,G)$ is the edge density of $G$.
For disjoint $A,B\subseteq V(G)$ we write $e(A,B)$ for the number of edges with
one end in $A$ and the other in $B$, and $\Delta(G)$ for the maximum degree of
$G$.  For a bipartition $V(G)=X\sqcup Y$, its internal-edge count is
$e(X)+e(Y)$, and $D_2(G)$ denotes the minimum of this quantity.

The max-cut theorem quoted in Section~\ref{sec:low} is proved only for
sufficiently large $n$.  Every normalized quantity used in this paper is
invariant under balanced blow-ups, and it is the last identity of the next
lemma that removes the restriction.
Write $G(k)$ for the balanced blow-up of
$G$ in which every vertex is replaced by $k$ pairwise nonadjacent twins.

\begin{lemma}\label{lem:blowup}
For every graph $G$ on $n$ vertices and every $k\ge1$,
\[
 \optf(G(k))=\optf(G),\qquad
 \frac{2e(G(k))}{(nk)^{2}}=\frac{2e(G)}{n^{2}},\qquad
 \frac{\hom(C_4,G(k))}{(nk)^{4}}
 =\frac{\hom(C_4,G)}{n^{4}},
\]
and
\[
 \frac{D_2(G(k))}{(nk)^{2}}=\frac{D_2(G)}{n^{2}}.
\]
\end{lemma}

\begin{proof}
Twin classes are independent, so an edge of $G(k)$ joins two classes that are
adjacent in $G$.  Given $f:V(G(k))\to[0,1]$, let $x_i$ be the mean of $f$ on
class $i$.  Then
$\sum_{\{u,v\}\in E(G(k))}f_uf_v=k^{2}\sum_{\{i,j\}\in E(G)}x_ix_j$, while the mass
constraint $\sum_uf_u=nk/2$ reads $\sum_ix_i=n/2$.  Dividing by $(nk)^{2}$
identifies the two minimization problems, so $\optf(G(k))=\optf(G)$.  The edge and
homomorphism counts satisfy $e(G(k))=k^{2}e(G)$ and
$\hom(C_4,G(k))=k^{4}\hom(C_4,G)$, since a
homomorphism of the blow-up is a homomorphism of $G$ together with a choice of
twin at each of the four positions.

For the last identity, a bipartition of $G(k)$ puts a fraction $x_i\in[0,1]$
of class $i$ on the first side, and its internal-edge count is
\[
 k^{2}\sum_{\{i,j\}\in E(G)}\bigl(x_ix_j+(1-x_i)(1-x_j)\bigr),
\]
again because the classes are independent.  This function is affine in each
$x_i$ separately, hence attains its minimum over $[0,1]^{n}$ at a point of
$\{0,1\}^{n}$, that is, at a bipartition of $G$.  Therefore
$D_2(G(k))\ge k^{2}D_2(G)$.  Conversely every bipartition of $G$ lifts to one
of $G(k)$ with exactly $k^{2}$ times the internal-edge count, so
$D_2(G(k))=k^{2}D_2(G)$.
\end{proof}

\section{The high-density range}\label{sec:high}

\begin{proposition}\label{prop:c4}
Let
\[
 \gamma=\frac3{64}+\frac1{50\,000\,000\,000}.
\]
Every triangle-free graph satisfies
\begin{equation}\label{eq:c4}
 t(C_4,G)\ge\dens^{3}-\gamma\dens.
\end{equation}
\end{proposition}

The slack $\varepsilon=1/(5\cdot10^{10})$ in $\gamma$ is an artefact of the
method rather than of the statement: \eqref{eq:c4} is very likely true with
$\gamma=3/64$, but that inequality is tight at the Clebsch graph, so every
square occurring in a certificate for it would have to vanish there as well,
and no such certificate was found.
The verifier reports that the exact slack its own square coefficients require
is $1.9911\cdot10^{-11}<\varepsilon$, and nothing in this paper needs the
sharp constant.

The value $3/64$ is, on the other hand, optimal.  Indeed, the Clebsch graph has
$\dens=\tfrac5{16}$ and $t(C_4,G)=\tfrac{1040}{65536}=\tfrac{65}{4096}$, while
$\dens^{3}-\tfrac3{64}\dens=\tfrac{125}{4096}-\tfrac{60}{4096}
=\tfrac{65}{4096}$ at $\dens=\tfrac5{16}$, so \eqref{eq:c4} with
$\gamma=\tfrac3{64}$ holds there with equality and fails for every smaller
$\gamma$.

\begin{proof}
The proof is a flag-algebra certificate: a finite identity exhibiting the
difference between the two sides of \eqref{eq:c4} as a nonnegative rational
combination of averaged squares, so that checking finitely many rational
inequalities establishes \eqref{eq:c4} for every triangle-free graph at once.

It is enough to prove the graphon statement, because a graph and its step
graphon have the same homomorphism densities: for the step graphon $W$ of $G$,
$t(K_2,W)=\dens$ and $t(C_4,W)=t(C_4,G)$.  At level six, the expression
$t(C_4,W)-\dens^{3}+\gamma\dens$ has one rational coefficient for each of the
$38$ triangle-free unlabelled graphs on six vertices.  The ancillary
certificate file
\[
 \texttt{clebsch\_tangent\_certificate.json}
\]
gives $34$ positive rational multiples of averaged squares of flags.  The
exact verifier
\[
 \texttt{verify\_clebsch\_tangent\_certificate.py}
\]
checks that the coefficient of every one of the $38$ hosts dominates the sum
of the corresponding square coefficients.  It also checks the host list,
flag dimensions, and positivity of every multiplier.  All sign decisions are
made in rational arithmetic.  This proves \eqref{eq:c4} by the flag-algebra
method \cite{razborov-flag}.
\end{proof}

\begin{lemma}\label{lem:anchor}
If $G$ is triangle-free, $\dens>0$, and $\Delta(G)<n/2$, then
\begin{equation}\label{eq:anchor}
 \optf(G)\le\frac{\dens}8-\frac{t(C_4,G)}{4\dens}.
\end{equation}
\end{lemma}

This is Proposition~1.1 of \cite{razborov-halves}, stated there for $\opt$
rather than for the relaxation $\optf$, and with the four-cycle density
normalized so that its value is $3\,t(C_4,G)$.  Razborov derives it from the
argument of \cite{krivelevich}.  We give the proof because the fractional form
is what the next section perturbs.

\begin{proof}
Fix an edge $\{u,v\}$ and put
\[
 A=\Nb(u),\qquad C=\Nb(v),\qquad
 B=V(G)\setminus(A\cup C).
\]
Triangle-freeness makes $A$ and $C$ disjoint independent sets.  Write
\[
 \delta_u=\frac{|A|}{n},\qquad \delta_v=\frac{|C|}{n},\qquad
 z=1-\delta_u-\delta_v,
\]
and
\[
 p=\frac{1/2-\delta_u}{z},\qquad
 1-p=\frac{1/2-\delta_v}{z}.
\]
The maximum-degree hypothesis gives $p\in[0,1]$.  The two functions
$\mathbf 1_A+p\mathbf 1_B$ and
$\mathbf 1_C+(1-p)\mathbf 1_B$ are feasible fractional halves.  If their
costs are $Q_u,Q_v$, then direct expansion gives
\begin{equation}\label{eq:edge-bound}
 (1-p)Q_u+pQ_v
 =p(1-p)\left(\frac{\dens}2-
 \frac{e(\Nb(u),\Nb(v))}{n^2}\right)
 \le\frac14\left(\frac{\dens}2-
 \frac{e(\Nb(u),\Nb(v))}{n^2}\right).
\end{equation}
Thus $\optf(G)$ is at most the right-hand side of \eqref{eq:edge-bound} for
every edge $\{u,v\}$.
Moreover,
\begin{equation}\label{eq:c4-count}
 \sum_{\{u,v\}\in E(G)}e(\Nb(u),\Nb(v))
 =\frac12\hom(C_4,G).
\end{equation}
Indeed, $\Nb(u)$ and $\Nb(v)$ are disjoint, so an edge counted by
$e(\Nb(u),\Nb(v))$ has exactly one end in $\Nb(u)$ and exactly one in
$\Nb(v)$.  Writing that end first,
\[
 e(\Nb(u),\Nb(v))=\#\{(x,y):x\in\Nb(u),\ y\in\Nb(v),\ \{x,y\}\in E(G)\}.
\]
Summing over the two orderings of each edge $\{u,v\}$ therefore gives
\[
 2\sum_{\{u,v\}\in E(G)}e(\Nb(u),\Nb(v))
 =\#\{(u,x,y,v):\{u,v\},\{u,x\},\{x,y\},\{y,v\}\in E(G)\},
\]
and the right side is exactly the number of homomorphisms $C_4\to G$, a
homomorphism being a cyclically adjacent quadruple.

The right-hand side of \eqref{eq:edge-bound} bounds $\optf(G)$ for each of the
$e(G)$ edges, hence so does its average, and by \eqref{eq:c4-count} that
average is
\[
 \frac1{e(G)}\sum_{\{u,v\}\in E(G)}\frac{e(\Nb(u),\Nb(v))}{n^{2}}
 =\frac{\hom(C_4,G)}{2e(G)n^{2}}
 =\frac{t(C_4,G)}{\dens},
\]
since $e(G)=\tfrac12\dens n^{2}$ and $\hom(C_4,G)=t(C_4,G)n^{4}$.  Therefore
$\optf(G)\le\tfrac14\bigl(\tfrac\dens2-t(C_4,G)/\dens\bigr)$, which is
\eqref{eq:anchor}.
\end{proof}

\begin{proposition}\label{prop:high}
If $G$ is triangle-free and $\dens\ge7/22$, then $\optf(G)<131/5000$.
\end{proposition}

\begin{proof}
If $\Delta(G)\ge n/2$, a neighbourhood of a maximum-degree vertex is an
independent set of size at least $n/2$.  Assigning values in $[0,1]$ on that
neighbourhood with total mass $n/2$, and zero elsewhere, gives a fractional
half of cost zero.  Thus suppose $\Delta(G)<n/2$.  Lemma~\ref{lem:anchor} and
Proposition~\ref{prop:c4} give
\[
 \optf(G)\le\frac{\dens}8-\frac{\dens^{2}}{4}+\frac\gamma4.
\]
This expression decreases for $\dens>1/4$.  Since $\gamma<293/6250$, at
$\dens=7/22$ it is less than
\[
 \frac7{484}+\frac{293}{25000}
 =\frac{79203}{3025000}<\frac{79255}{3025000}=\frac{131}{5000},
\]
where $\frac7{484}=\frac7{176}-\frac{49}{1936}$ is the value of
$\frac\dens8-\frac{\dens^{2}}4$ at $\dens=\frac7{22}$.
\end{proof}

\section{The low-density range}\label{sec:low}

\begin{theorem}[Balogh--Clemen--Lidick\'y \cite{bcl}]\label{thm:bcl-large}
There is an $n_0$ such that every triangle-free graph $G$ on $n\ge n_0$
vertices satisfies
\[
 D_2(G)\le\frac{n^{2}}{23.5}=\frac2{47}n^{2}.
\]
\end{theorem}

This is Theorem~2(a) of \cite{bcl}, where it is stated with
the constant $23.5$ and for $n$ large enough.  The largeness hypothesis is not
a restriction here.

\begin{corollary}\label{cor:bcl}
Every triangle-free graph $G$ on $n\ge1$ vertices satisfies
$D_2(G)\le\frac2{47}n^{2}$.
\end{corollary}

\begin{proof}
The blow-up $G(k)$ is triangle-free and has $nk$ vertices, so
Theorem~\ref{thm:bcl-large} applies to it once $k\ge n_0/n$.  By
Lemma~\ref{lem:blowup}, $D_2(G)/n^{2}=D_2(G(k))/(nk)^{2}\le2/47$.
\end{proof}

\subsection{A perturbation inside one neighbourhood}
A weighted graph is a finite graph $H$ with positive vertex weights
$(w_x)$ summing to one.  By a \emph{profile} on $H$ we mean a function
$\pi:V(H)\to[0,1]$.  Its weighted mean is $\sum_xw_x\pi_x$.  A constant profile
is one that is constant on $V(H)$.  Put
\[
 \mu=\sum_{\{x,y\}\in E(H)}w_xw_y,\qquad
 d_x=\sum_{y\in\Nb(x)}w_y,\qquad
 e_x=\sum_{y\in\Nb(x)}w_yd_y,
\]
and
\[
 S_x=d_x(e_x-\mu d_x).
\]
We call $\mu$ the \emph{edge mass} of $H$.

\begin{lemma}\label{lem:cubic}
If $\mu>0$, some positive-degree vertex $x$ satisfies
\begin{equation}\label{eq:cubic}
 S_x\ge4\mu^3.
\end{equation}
\end{lemma}

\begin{proof}
Put $p=2\mu$ and, for $d_x>0$, define
\[
 R_x=\frac{d_x}{p},\qquad
 U_x=\frac1{d_x}\sum_{y\in\Nb(x)}w_yR_y.
\]
Then $\sum_xw_xR_x=1$, $e_x=p^2R_xU_x$, and hence
\begin{equation}\label{eq:S-RU}
 S_x=p^3R_x^2\left(U_x-\frac12\right).
\end{equation}
Suppose that \eqref{eq:cubic} fails strictly at every nonisolated vertex.
Since $4\mu^3=p^3/2$, equation \eqref{eq:S-RU} gives
\[
 U_x<\frac{R_x^2+1}{2R_x^2}.
\]
The weighted harmonic--arithmetic mean inequality in $\Nb(x)$ therefore
gives
\begin{equation}\label{eq:harmonic}
 \frac1{d_x}\sum_{y\in\Nb(x)}\frac{w_y}{R_y}
 \ge\frac1{U_x}>
 \frac{2R_x^2}{R_x^2+1}.
\end{equation}
No neighbour in this expression is isolated.

Let $\sigma$ be the total weight of the nonisolated vertices.  Average
\eqref{eq:harmonic} with weights $w_xR_x$, which sum to one.  The average of
the left side is exactly
\[
 \frac1p\sum_{y:d_y>0}\frac{w_y}{R_y}
 \sum_{x\in\Nb(y)}w_x
 =\frac1p\sum_{y:d_y>0}\frac{w_yd_y}{R_y}=\sigma.
\]
For every $R\ge0$,
\[
 \frac{2R^3}{R^2+1}-(2R-1)
 =\frac{(R-1)^2}{R^2+1}\ge0.
\]
Consequently the average of the right side of \eqref{eq:harmonic} is at
least
\[
 \sum_{x:d_x>0}w_x(2R_x-1)=2-\sigma\ge\sigma,
\]
contradicting the strict inequality in \eqref{eq:harmonic}.
\end{proof}

For $\alpha\in[0,1]$ define
\[
 \hat\alpha=\min(\alpha,1-\alpha),\qquad
 L(\alpha)=\min\{4\hat\alpha^{2},\hat\alpha(1-\hat\alpha)\}.
\]

\begin{lemma}\label{lem:onestep}
Let $H$ be a weighted triangle-free graph of edge mass $\mu$, and let
$\psi:V(H)\to\mathbb R$ be arbitrary.  For
\[
 \mathcal A(\pi)=\sum_{\{u,v\}\in E(H)}w_uw_v\pi_u\pi_v
 +\sum_uw_u\psi_u\pi_u,
\]
the constant profile $\alpha$ can be replaced by a profile
$\pi:V(H)\to[0,1]$ with the same weighted mean and
\begin{equation}\label{eq:onestep}
 \mathcal A(\pi)\le\mathcal A(\alpha)-4L(\alpha)\mu^3.
\end{equation}
\end{lemma}

\begin{proof}
The assertion is immediate if $\mu=0$ or $\alpha\in\{0,1\}$.  Otherwise
choose $x$ from Lemma~\ref{lem:cubic}, abbreviate
$d=d_x$, $z=1-d$, and put
\[
 \varphi=\mathbf 1_{\Nb(x)}-d.
\]
This function has weighted mean zero.  Since $H$ is triangle-free,
$\Nb(x)$ is independent, and expanding the quadratic form gives
\begin{equation}\label{eq:quad}
 -\sum_{\{u,v\}\in E(H)}w_uw_v\varphi_u\varphi_v
 =d(e_x-\mu d)=S_x.
\end{equation}
The positive vertex weights and $d_x>0$ give $0<d,z<1$, so all radii below
are well defined.

The full feasible interval for $\alpha+t\varphi$ is
$-t_-\le t\le t_+$, where
\[
 t_+=\min\left\{\frac{1-\alpha}{z},\frac\alpha d\right\},
 \qquad
 t_-=\min\left\{\frac\alpha z,\frac{1-\alpha}{d}\right\}.
\]
Give the two endpoints the barycentric weights
\[
 \lambda_+=\frac{t_-}{t_++t_-},\qquad
 \lambda_-=\frac{t_+}{t_++t_-}.
\]
The linear terms in $t$ cancel, including the arbitrary $\psi$-term, while
\eqref{eq:quad} gives
\[
 \lambda_+\mathcal A(\alpha+t_+\varphi)
 +\lambda_-\mathcal A(\alpha-t_-\varphi)
 =\mathcal A(\alpha)-t_+t_-S_x.
\]
At least one endpoint is no larger than this average.

It remains to bound the product of the radii.  By symmetry assume
$\alpha\le1/2$.  Directly from their definitions,
\[
 t_+t_-=\begin{cases}
 \alpha(1-\alpha)/z^2,&d\le\alpha,\\[3pt]
 \alpha^2/(dz),&\alpha\le d\le1-\alpha,\\[3pt]
 \alpha(1-\alpha)/d^2,&d\ge1-\alpha.
 \end{cases}
\]
The first and third quantities are at least $\alpha(1-\alpha)$, and the
middle one is at least $4\alpha^2$.  Hence $t_+t_-\ge L(\alpha)$.
Together with Lemma~\ref{lem:cubic}, this proves \eqref{eq:onestep}.
\end{proof}

\subsection{The cut construction}\label{ssec:low-cut}
Assume from now on that
\begin{equation}\label{eq:low-density}
 \dens\le\dens_0:=\frac7{22}.
\end{equation}
By Corollary~\ref{cor:bcl}, choose a bipartition $V(G)=X\sqcup Y$ with
\[
 I:=\frac{e(X)+e(Y)}{n^2}\le I_0:=\frac2{47}.
\]
Interchanging the sides if necessary, write
\[
 \frac{|X|}{n}=\frac{1-s}{2},\qquad
 \frac{|Y|}{n}=\frac{1+s}{2},\qquad 0\le s\le1,
\]
and put
\[
 a=\frac{s}{1+s}\in\left[0,\frac12\right],
 \qquad c=1-a,\qquad
 I_Y=\frac{e(Y)}{n^2},\qquad
 \mu=4c^2I_Y.
\]
The graph $G[Y]$, with every vertex assigned weight $1/|Y|$, has edge mass
$\mu$, because $|Y|/n=1/(2c)$.  Also
\begin{equation}\label{eq:mu-cap}
 0\le\mu\le4I_0c^2.
\end{equation}

We define two fractional halves:
\[
 \mathbf 1_X+a\mathbf 1_Y,
 \qquad c\mathbf 1_Y.
\]
Their respective costs are
\begin{equation}\label{eq:baselines}
 B_1=\frac{a\dens}{2}+cI-\frac{1+a}{4c}\mu,
 \qquad B_2=\frac\mu4.
\end{equation}
Indeed, if $I_X=e(X)/n^2=I-I_Y$, then the first cost is
$I_X+a(\tfrac{\dens}2-I)+a^2I_Y$, and the second is $c^2I_Y$.

Apply Lemma~\ref{lem:onestep} inside $Y$.  Edges from $Y$ to the fixed
profile on $X$ contribute only a linear functional, so the lemma applies to
both baselines.  Their constant values on $Y$ are $a$ and $c$, and
$L(a)=L(c)$.  Returning from the normalization of $G[Y]$ to that of $G$
multiplies the saving by $(|Y|/n)^2=1/(4c^2)$.  Preservation of the weighted
mean on $Y$ keeps the resulting ambient profile a fractional half.  We have
proved
\begin{equation}\label{eq:master}
 \optf(G)\le\min(B_1,B_2)-\frac{L(a)\mu^3}{c^2}.
\end{equation}

Increasing $\dens$ and $I$ can only increase $B_1$.  Thus it remains to
maximize the following elementary expression with $\dens=\dens_0$ and $I=I_0$:
\begin{equation}\label{eq:F-two}
 F(a,\mu)=
 \min\left\{P(a)-\frac{1+a}{4c}\mu,\frac\mu4\right\}
 -\frac{L(a)}{c^2}\mu^3,
 \qquad
 P(a)=\frac{a\dens_0}{2}+cI_0,
\end{equation}
on $0\le a\le1/2$ and the interval \eqref{eq:mu-cap}.

\subsection{The one-variable calculation}
\begin{lemma}\label{lem:calculus}
On the domain just described,
\[
 F(a,\mu)<\frac{131}{5000}.
\]
\end{lemma}

\begin{proof}
This is a straightforward calculation: we maximize in $\mu$ for each fixed $a$,
which leaves a one-variable envelope, and then locate the maximum of that
envelope using coarse rational bounds on its derivative.

For fixed $a$, the two quantities inside the minimum in \eqref{eq:F-two} are
equal precisely when
\[
 \mu=\nu(a):=2cP(a),
\]
the first exceeding the second for $\mu>\nu(a)$ and the second exceeding the
first for $\mu<\nu(a)$.  For $\mu\le\nu(a)$, where the minimum equals
$\tfrac\mu4$, the derivative of \eqref{eq:F-two} with respect to $\mu$ is
\[
 \frac14-\frac{3L(a)\mu^2}{c^2}
 \ge\frac14-12I_0^2
 =\frac14-\frac{48}{2209}>0,
\]
where we used $L(a)\le1/4$ and \eqref{eq:mu-cap}.  For $\mu\ge\nu(a)$, where
the minimum equals $P(a)-\tfrac{(1+a)\mu}{4c}$, that derivative is strictly
negative.  Hence the maximizing value of $\mu$ is
\[
 \min\{\nu(a),4I_0c^2\}.
\]
The two entries agree exactly when
\[
 a=a_*:=\frac{I_0}{\dens_0/2+I_0}=\frac{88}{417}.
\]

Put $\delta=\tfrac{\dens_0}2-I_0=241/2068$.  Since
\[
 L(a)=\begin{cases}4a^2,&a\le1/5,\\a(1-a),&a\ge1/5,
 \end{cases}
\]
maximizing first in $\mu$ gives the one-variable envelope
\begin{equation}\label{eq:envelope}
 \begin{cases}
 \displaystyle \frac{cP}{2}-32a^2cP^3,
   &0\le a\le\frac15,\\[5pt]
 \displaystyle \frac{cP}{2}-8ac^2P^3,
   &\frac15\le a\le a_*,\\[5pt]
 \displaystyle I_0c^2-64I_0^3ac^5,
   &a_*\le a\le\frac12.
 \end{cases}
\end{equation}

We check its monotonicity with deliberately coarse rational bounds.  Put
$h=cP/2$.  On $[0,a_*]$,
\[
 h'=\frac{\dens_0/2-2I_0-2\delta a}{2}
 \ge h'(a_*)=\frac{455}{36696}>\frac1{100}.
\]
Also $P\le P(a_*)=28/417<7/100$ and $\delta<1/8$.

On $0\le a\le1/5$, the derivative of the subtracted term in the first line
of \eqref{eq:envelope} is
\[
 32aP^2\bigl((2-3a)P+3ac\delta\bigr)
 <32\cdot\frac15\left(\frac7{100}\right)^2
 \left(2\cdot\frac7{100}+\frac3{40}\right)
 <\frac7{1000}<\frac1{100}.
\]
Thus the first line is increasing.  On $1/5\le a\le a_*$ we have
$a_*<17/80$ and $c\le4/5$.  The derivative of the subtracted term in the
second line is
\[
 8cP^2\bigl((1-3a)P+3ac\delta\bigr)
 <8\cdot\frac45\left(\frac7{100}\right)^2
 \left(\frac25\cdot\frac7{100}
 +3\cdot\frac{17}{80}\cdot\frac45\cdot\frac18\right)
 <\frac1{300}<\frac1{100}.
\]
The second line is increasing as well.

For the third line of \eqref{eq:envelope}, differentiation gives
\[
 -2I_0c\left[1-32I_0^2c^3(6a-1)\right].
\]
The bracket is positive, since
\[
 32I_0^2c^3(6a-1)\le64I_0^2=\frac{256}{2209}<1.
\]
Thus the third line is decreasing.  The global maximum is attained at
$a_*=88/417$.

At that point
\[
 c_*=\frac{329}{417},\qquad
 \mu_*=4I_0c_*^2=\frac{18424}{173889},\qquad
 \frac{a_*}{c_*}=\frac{88}{329}.
\]
The baseline term satisfies
\[
 I_0c_*^2=\frac{4606}{173889}<\frac{2649}{100000},
\]
and the saving satisfies
\[
 \frac{a_*}{c_*}\mu_*^3
 >\frac4{15}\left(\frac{21}{200}\right)^3
 >\frac{29}{100000}.
\]
The two lower bounds used here are
\[
 \frac{88}{329}>\frac4{15},\qquad
 \frac{18424}{173889}>\frac{21}{200},
\]
and every comparison is a direct cross-multiplication.  Therefore
\[
 F(a_*,\mu_*)
 <\frac{2649}{100000}-\frac{29}{100000}
 =\frac{131}{5000},
\]
as required.
\end{proof}

\begin{proof}[Proof of Theorem~\ref{thm:main}]
By Lemma~\ref{lem:rounding}, it is enough to bound $\optf(G)$.  If
$\dens\ge7/22$,
apply Proposition~\ref{prop:high}.  Otherwise the cut construction of
Section~\ref{ssec:low-cut} and Lemma~\ref{lem:calculus} give
$\optf(G)<131/5000$.
\end{proof}

\section{Verification}\label{sec:verification}

The proof supplies one machine-certified component,
Proposition~\ref{prop:c4}.  With Python and NetworkX installed, the ancillary
files can be checked from scratch with
\[
 \texttt{python3 verify\_clebsch\_tangent\_certificate.py}.
\]
The verifier works in exact rational arithmetic, so no floating-point
comparison enters any sign decision.  Apart from this certificate and the
max-cut theorem quoted as Theorem~\ref{thm:bcl-large}, every step of the
argument is carried out by hand in the text above.

The argument has also been formalized in the Lean~4 proof assistant, on top of
\texttt{mathlib}.  The ancillary directory \texttt{lean/} contains the
development, and \texttt{./verify.sh} builds it and prints an audit of the
axioms each conclusion depends on.

One statement is not formalized: the max-cut bound of
Theorem~\ref{thm:bcl-large} is taken as an explicit hypothesis of the
formalized theorem, so it appears in
its statement rather than being reproved.  Everything else is proved inside the
development, including Proposition~\ref{prop:c4} together with the
flag-algebra machinery it needs: the sampling semantics, the exhaustive check
over the triangle-free six-vertex hosts, and the passage from the certificate
to the graph-level inequality.

There is no use of \texttt{sorry}, \texttt{native\_decide}, or any added
axiom, and each conclusion depends on
exactly the three axioms of Lean's standard library, namely \texttt{propext},
\texttt{Classical.choice} and \texttt{Quot.sound}.

\begin{thebibliography}{99}

\bibitem{bcl}
J.~Balogh, F.~C.~Clemen and B.~Lidick\'y,
\emph{Max cuts in triangle-free graphs},
in: Extended Abstracts EuroComb 2021, Trends in Mathematics, Springer, 2021,
509--514; \texttt{arXiv:2103.14179v1}.

\bibitem{erdosproblems}
T.~F.~Bloom,
\emph{Erd\H{o}s Problem \#128},
\url{https://www.erdosproblems.com/128}, accessed 2026-08-03.

\bibitem{erdos75}
P.~Erd\H{o}s,
\emph{Problems and results in graph theory and combinatorial analysis},
in: Proc. Fifth British Combinatorial Conference (Univ. Aberdeen, 1975),
Congr. Numer. 15 (1976), 169--192.

\bibitem{erdos97}
P.~Erd\H{o}s,
\emph{Some old and new problems in various branches of combinatorics},
Discrete Math. 165/166 (1997), 227--231; proceedings of the international
symposium Graphs and Combinatorics (Marseille, 1995).

\bibitem{efps}
P.~Erd\H{o}s, R.~J.~Faudree, J.~Pach and J.~Spencer,
\emph{How to make a graph bipartite},
J. Combin. Theory Ser. B 45(1) (1988), 86--98.

\bibitem{efrs}
P.~Erd\H{o}s, R.~J.~Faudree, C.~C.~Rousseau and R.~H.~Schelp,
\emph{A local density condition for triangles},
Discrete Math. 127(1--3) (1994), 153--161.

\bibitem{egs}
P.~Erd\H{o}s, E.~Gy\H{o}ri and M.~Simonovits,
\emph{How many edges should be deleted to make a triangle-free graph
bipartite?},
in: Sets, graphs and numbers (Budapest, 1991), Colloq. Math. Soc. J\'anos
Bolyai 60, North-Holland, Amsterdam, 1992, 239--263.

\bibitem{keevash-sudakov}
P.~Keevash and B.~Sudakov,
\emph{Sparse halves in triangle-free graphs},
J. Combin. Theory Ser. B 96(4) (2006), 614--620.

\bibitem{krivelevich}
M.~Krivelevich,
\emph{On the edge distribution in triangle-free graphs},
J. Combin. Theory Ser. B 63(2) (1995), 245--260.

\bibitem{norin-yepremyan}
S.~Norin and L.~Yepremyan,
\emph{Sparse halves in dense triangle-free graphs},
J. Combin. Theory Ser. B 115 (2015), 1--25.

\bibitem{razborov-flag}
A.~A.~Razborov,
\emph{Flag algebras},
J. Symbolic Logic 72(4) (2007), 1239--1282.

\bibitem{razborov-halves}
A.~A.~Razborov,
\emph{More about sparse halves in triangle-free graphs},
Sb. Math. 213(1) (2022), 109--128.

\end{thebibliography}

\end{document}
